\scanpage{235}\mspage{90}
7. \uline{Morphismes cofiltrants et loc cofiltrants}.

Mais je ne veux pas m'attarder à présent sur
ces développements, qui ne me paraissent pas
urgents, et je voudrais revenir aux questions
liées à la notion de catégories filtrantes \deleted{\unsure{et}}
cofiltrantes.

% slanted
\margin{J'ai \unsure{modifié} la terminologie, \unsure{maintenant} \unsure{cof}, loc. fbt cof, \unsure{pré} \illegible{1} ; fbt cof = \unsure{$0$-connexité|$0$-conn.} \unsure{stable par passage|stable par prod. fibrés} : \unsure{si|tt} $X'\unsure{-|\to}\unsure{X|x}$ \unsure{dans} $X$ cof, $X'$ étale et \unsure{connexe}.}
% slanted
\margin{Revoir la terminologie du cas (loc. fbt cofiltrant, \unsure{etc.}), qui est une situation très \unsure{ingrate}, faute d'être stable par passage \unsure{aux|un} $X^{\unsure{\prime|\prime\prime}}$ \unsure{étal|étale} sur $X$.}
\uline{Définition} 7.1. Soit $f:X\to Y$ une flèche dans $\operatorname{Cat}$.
On dit que $f$ \uline{est cofiltrant} (resp. \uline{loc. cofiltrant},
\uline{faiblement cofiltrant}$^{*}$, respt \uline{loc. faiblement cofiltrant}
-- ou \uline{précofiltrant}) si pour tt $y$ dans
$Y$, la catégorie \inserted{colocalisée} $\,{}_{y\backslash}X$ est cofiltrante (resp. \ldots)

[\inserted{Dualmt,} On dit que $f$ est filtrant, resp. \ldots, si les
catégories $X_{/y}$ sont filtrantes, resp. \ldots, i.e.
si $f^\circ:X^\circ\unsure{-|\to}Y^\circ$ est cofiltrant resp \ldots]

\deleted{\unsure{Proposition|Remarque}} \uline{\unsure{Remarque} 7.2.} \unsure{Notions} stables par chgt de base coétale
\inserted{contient les iso, \unsure{et|de} $=$ \unsure{bien sûr|les quatre genres} \unsure{imm}}\unsure{ersions}
% slanted
\margin{\deleted{\illegible{1}} et \unsure{les} \unsure{immersions} \unsure{coétales} sont \unsure{plates}.}
\uline{Proposition} 7.3. Soit $f:X\unsure{-|\to}Y$ une \deleted{flèche}
cofiltrant (resp. loc. cofiltrant, \unsure{respt.} fbt cofiltrant,
resp. \deleted{loc. fbt \unsure{pr}} précofiltrant), et supposons
que $Y$ soit cofiltrant (resp. loc cofiltrant,
resp fbt cofiltrant, resp. précofiltrant). Alors
$X$ lui-même est cofiltrant (resp. \ldots).

\uline{Dém}. Il s'agit, dans les quatre cas
envisagés, de vérifier pour $X$ la condition
\[
\underbrace{\mathrm{PS}_g1'\ /}_{\text{\unsure{deux autres cas|deux premiers cas}}}
\ \mathrm{PS}_g1\ /\ \mathrm{PS}_g1'+\mathrm{PS}_g2\ /\ \mathrm{PS}_g1+\mathrm{PS}_{\unsure{g|s}}2.
\]
Mais
% archive sheet number: 235
